%!TeX root=MemoriaTFG.tex

\chapter{Discusión}
\label{cap:discusion}

\section{Significado de los resultados}
\label{sec:significado}

\subsection{Especialización del dominio frente a encoders generalistas}
\label{sec:disc-especializacion}

Los resultados del sondeo lineal (secci\'on~\ref{sec:rep-panderm} y
secci\'on~\ref{sec:benchmark}) sostienen la hipótesis de que el
preentrenamiento sobre datos del dominio clínico aporta una ventaja
empírica medible sobre encoders generalistas adaptados. PanDerm
Large supera a PanDerm Base en accuracy, AUROC y F1 ponderada sobre
los diez datasets externos contrastados, con mejora media de
$+4{,}1$\,pp en accuracy, $+9{,}0$\,pp en BAcc y
$+3{,}6$\,pp en AUROC. La mejora absoluta es máxima sobre
Dermnet ($+10{,}0$\,pp accuracy), WSI~patches
($+8{,}7$\,pp) y PAD-UFES-20 ($+4{,}7$\,pp), datasets con
mayor número de clases o mayor heterogeneidad de adquisición.

La comparación entre PanDerm Large y la mejor arquitectura
convolucional supervisada cuantifica la magnitud del beneficio del
preentrenamiento de dominio en función de la dificultad del
dataset. Sobre HAM10000 (dermatoscopia estandarizada), la brecha
en accuracy es modesta ($+0{,}5$\,pp) y ConvNeXt-Large incluso
supera a PanDerm Large en AUROC ($0{,}967$ vs $0{,}954$,
$-1{,}3$\,pp), comportamiento compatible con el hecho de que
las características generales de ImageNet-21K capturan ya patrones
suficientemente informativos para el ranking de probabilidades en
una modalidad estandarizada. Sobre PAD-UFES-20 (fotografía clínica
con dispositivo móvil), la brecha es de $+8{,}2$\,pp accuracy y
$+5{,}4$\,pp AUROC; sobre DDI (fotografía clínica con
representación equilibrada de fototipos), $+2{,}9$\,pp accuracy
y $+8{,}6$\,pp AUROC. El patrón sugiere que la ventaja del
preentrenamiento de dominio crece con la distancia entre el dataset
y los corpus web generalistas, y resulta especialmente relevante en
condiciones clínicas no estandarizadas (atención primaria,
teledermatología con dispositivos heterogéneos).

SigLIP-Large SO400M, modelo vision-lenguaje generalista, obtiene
la mejor accuracy y AUROC sobre HAM10000 ($0{,}900 / 0{,}971$),
comportamiento atribuible a la sobrerrepresentación de dermatoscopia
estandarizada en los $400$ millones de pares imagen-texto web de
su preentrenamiento. La superioridad de SigLIP-Large desaparece
sobre PAD-UFES-20 y DDI, lo que delimita su rango de utilidad: es
competitivo en modalidades canónicas pero no en material clínico
heterogéneo.

La caracterización del preentrenamiento específico de dominio
admite además una segunda lectura sobre el nivel ontológico de
subcategoría clínica. El experimento de combinación de
codificadores sobre DermapixelAI~1.0
(secci\'on~\ref{sec:dermapixel-ensemble},
tabla~\ref{tab:dermapixel-ensemble}) sitúa a DermLIP v2 como la
mejor cabeza individual sobre L2 con AUROC $0{,}860$,
$+6{,}7$\,pp por encima de PanDerm Large ($0{,}793$) y
$+15{,}3$\,pp por encima de PanDerm Base ($0{,}707$).
El hallazgo es coherente con la lectura \emph{zero-shot} sobre el
mismo corpus (secci\'on~\ref{sec:dermapixel-zs}), donde DermLIP v2
lidera también la AUROC L2 sin entrenar cabeza alguna
($0{,}794$ en inglés). La interpretación conjunta sugiere que el
preentrenamiento contrastivo texto-imagen sobre el millón de pares
de Derm1M~\cite{derm1m} captura la granularidad de subcategoría
clínica con mayor fidelidad que el preentrenamiento
auto-supervisado puro de la familia PanDerm, lo que matiza la
jerarquía PanDerm Large $>$ PanDerm Base $>$ DermLIP v2 observada
en la mayoría de datasets y niveles: la elección del codificador
óptimo depende del nivel ontológico, y el preentrenamiento
multimodal aporta una señal complementaria sobre la categorización
clínica intermedia que no se reduce al beneficio de las
\emph{features} visuales aisladas.

El experimento de ajuste fino parcial sobre DermapixelAI~1.0 nivel
L1 (secci\'on~\ref{sec:dermapixel-ft-l1},
tabla~\ref{tab:dermapixel-ft-final}) refuerza esta lectura desde
un ángulo metodológico. Con $N_{\mathrm{train}} = 874$ y el
encoder PanDerm Large parcialmente descongelado
($25{,}2$\,M parámetros entrenables), el FT mejora la BAcc en
$+18{,}2$\,pp sobre la cifra robusta del sondeo lineal congelado
estimada por validación cruzada agrupada por caso
($0{,}440 \to 0{,}622$), pero no supera al \emph{ensemble} lineal
\texttt{avg\_large\_dermlip} de la
secci\'on~\ref{sec:dermapixel-ensemble} (BAcc $0{,}622$ vs
$0{,}702$, AUROC $0{,}715$ vs $0{,}817$). El patrón es coherente
con la literatura sobre régimen de bajo $N$: una combinación
lineal de codificadores fundacionales congelados aprovecha la
diversidad de sus representaciones pre-entrenadas sin introducir
los parámetros adicionales que el FT denso requiere ajustar a
partir de un corpus reducido. La implicación de diseño para
SpanDerm v0 (secci\'on~\ref{sec:linea-spanderm}) es directa: las
líneas prioritarias son la adaptación de bajo rango mediante
LoRA y la combinación de codificadores con cabezas ligeras antes
que el FT denso clásico, opción metodológica que las cifras
empíricas sobre el corpus en castellano justifican. La ejecución
posterior de SpanDerm v0
(secci\'on~\ref{sec:dermapixel-spanderm-v0}) confirma
cuantitativamente esta implicación: la adaptación LoRA sobre los
dos últimos bloques de PanDerm Large con sólo $524\,K$
parámetros entrenables ($0{,}17\%$ del total) alcanza BAcc
$= 0{,}363 \pm 0{,}007$ sobre L2, $+19{,}5$\,pp sobre la cifra
de validación cruzada del sondeo lineal congelado y $+8{,}4$\,pp
sobre DermLIP v2 individual, situándose como la mejor BAcc L2
documentada sobre el corpus. La lectura cruzada con DermLIP v2 es
particularmente informativa: la AUROC L2 de SpanDerm v0
($0{,}855 \pm 0{,}006$) empata estadísticamente con la de
DermLIP v2 nativo ($0{,}860$), lo que sugiere que el techo
discriminativo del corpus a nivel L2 es alcanzable tanto por el
preentrenamiento contrastivo PubMedBERT-DermLIP nativo como por
la adaptación textual ligera vía LoRA de un encoder de mayor
capacidad. La elección del codificador base y la estrategia de
adaptación operan como dimensiones aproximadamente ortogonales
sobre la cardinalidad L2.

La extensión posterior del FT denso parcial a los niveles L2 y L3
(secci\'on~\ref{sec:dermapixel-ft},
tabla~\ref{tab:dermapixel-ft-multilevel}) confirma el patrón
observado en L1 sobre la cardinalidad clínica L2: el FT denso con
$25{,}2$\,M parámetros entrenables alcanza BAcc L2 $0{,}324$ y
AUROC L2 $0{,}852$, sin superar a SpanDerm v0 LoRA ($524\,K$
parámetros entrenables, BAcc $0{,}363$, AUROC $0{,}855$). La
diferencia de $-4{,}0$\,pp BAcc y empate en AUROC sobre un
encoder $48$ veces menos parametrizado refuerza la hipótesis
metodológica formulada en secci\'on~\ref{sec:disc-lp-vs-ft}: bajo
$N_{\mathrm{train}} = 874$, la adaptación de bajo rango sobre un
encoder fundacional congelado domina al FT denso clásico sobre el
mismo encoder. Sobre L3, el FT denso no mejora al sondeo lineal
($+0{,}0$\,pp BAcc, $-0{,}9$\,pp AUROC), lo que confirma que la
cola larga estructural ($3{,}9$ muestras por clase de media en
\emph{train}) constituye un techo absoluto sobre cualquier
protocolo de ajuste denso sobre el corpus actual.

La extensión a codificadores generalistas no médicos
(secci\'on~\ref{sec:dermapixel-extra-encoders}) refuerza la lectura
de no-uniformidad: CLIP-L-336 generalista, sin material biomédico
en su preentrenamiento, encabeza la AUROC sobre L2
($0{,}826$ vs $0{,}793$ PanDerm Large, $+3{,}3$\,pp) y sobre L3
($0{,}827$ vs $0{,}813$, $+1{,}4$\,pp) bajo sondeo lineal, lo que
matiza la hipótesis de especialización de dominio sobre las
métricas de ranking probabilístico. La sustitución del clasificador
paramétrico por \emph{retrieval} $k$-NN con índice FAISS sobre los
mismos \emph{embeddings} PanDerm Large
(secci\'on~\ref{sec:dermapixel-extra-faiss}) supera al sondeo
lineal sobre el régimen de cola larga severa L3
($+2{,}6$\,pp BAcc, $+3{,}5$\,pp Acc@1), evidencia adicional
de que ningún codificador ni clasificador único domina en todos
los niveles ontológicos y de que la elección óptima debe
condicionarse al nivel objetivo y al criterio operativo
predominante.

\subsection{Eficiencia en etiquetas y régimen de saturación}
\label{sec:disc-eficiencia}

La curva de eficiencia en etiquetas
(secci\'on~\ref{sec:label-eff-results}, tabla~\ref{tab:label-eff})
documenta un fenómeno consistente con encoders preentrenados sobre
corpus masivos del dominio: la saturación del rendimiento sobre
HAM10000 ocurre con un volumen de datos etiquetados muy reducido.
Con $1\%$ del conjunto de entrenamiento ($82$ imágenes),
PanDerm Large alcanza $0{,}850$ accuracy y $0{,}918$ AUROC, lo
que equivale al $95{,}7\%$ y al $96{,}2\%$ del rendimiento
con el conjunto completo respectivamente. Con $5\%$
($410$ imágenes), la AUROC sube a $0{,}932$, equivalente al
$97{,}7\%$ del rendimiento máximo; la diferencia con
$100\%$ es de sólo $2{,}2$\,pp. La BAcc, en cambio, escala
con mayor lentitud y alcanza su máximo en $50\%$ del conjunto
($0{,}589$), comportamiento consistente con la dependencia de las
clases minoritarias frente al volumen total.

Sobre PAD-UFES-20 ($N_{\mathrm{train}} = 1\,493$ con el conjunto
completo), la saturación es aún más temprana: con $14$ imágenes
($1\%$) se alcanza $0{,}740$ accuracy y $0{,}907$ AUROC, lo
que equivale al $95{,}5\%$ del AUROC con el conjunto completo. La
implicación clínica es directa: en contextos con disponibilidad
limitada de datos etiquetados (enfermedades raras, centros de bajo
volumen, subgrupos poco representados), una cabeza lineal sobre
PanDerm Large preentrenado constituye una opción técnicamente
defendible.

\subsection{Sondeo lineal frente a ajuste fino}
\label{sec:disc-lp-vs-ft}

La comparación directa entre sondeo lineal y ajuste fino sobre
cuatro datasets (secci\'on~\ref{sec:ft-results},
tabla~\ref{tab:ft-panderm}) identifica un patrón asimétrico cuya
interpretación operativa es relevante para el diseño de sistemas
clínicos. Sobre HAM10000 y PAD-UFES-20, datasets con tamaño de
entrenamiento intermedio o alto ($N \in [1\,493, 8\,207]$), el
ajuste fino mejora la BAcc en $+47{,}1$\,pp y $+19{,}4$\,pp
respectivamente, lo que confirma la utilidad de la receta original
de PanDerm con \emph{weighted random sampler}, Mixup, CutMix y
\emph{layer decay} para el aprendizaje de las clases minoritarias
en presencia de desbalance acusado.

Sobre DDI, con $N_{\mathrm{train}} = 427$, el comportamiento se
invierte: el sondeo lineal supera al ajuste fino en accuracy
($+7{,}3$\,pp), BAcc ($+2{,}1$\,pp), AUROC
($+14{,}5$\,pp) y F1 ponderada ($+5{,}6$\,pp). El descenso
de AUROC ($0{,}827 \to 0{,}682$) es particularmente diagnóstico:
indica que el encoder ajustado pierde capacidad de ranking sobre el
\emph{test}, comportamiento compatible con sobreajuste a la
distribución de entrenamiento pese a las augmentaciones
empleadas. Sobre MSKCC ($N_{\mathrm{train}} = 6\,189$), la mejora
de BAcc es marginal ($+3{,}0$\,pp) y la AUROC desciende
ligeramente ($-2{,}7$\,pp).

Estos tres casos delimitan un régimen operativo cuantificable:
para datasets con $N_{\mathrm{train}} \lesssim 500$ imágenes,
el sondeo lineal sobre PanDerm Large constituye la opción de
referencia; para $N_{\mathrm{train}} \gtrsim 5\,000$, el ajuste
fino con la receta original es preferible siempre que existan
clases minoritarias clínicamente relevantes que el sondeo lineal
no logre discriminar. Entre ambos regímenes, la decisión depende
del desbalance de clases y del coste computacional asumible.

\subsection{Segmentación promptable y adaptación de bajo coste}
\label{sec:disc-segmentacion}

Los resultados de segmentación binaria sobre ISIC2018
(secci\'on~\ref{sec:seg-results}, tabla~\ref{tab:seg}) son el hallazgo
con mejor magnitud absoluta del trabajo: SAM2.1-Large con
adaptación LoRA ($r=8$) alcanza Dice $0{,}947$ sobre el conjunto
de test de ISIC2018, superando en $+2{,}6$\,pp la cifra
publicada por Yan et~al.~\cite{panderm2025} ($0{,}921$) y en
$+5{,}3$\,pp al baseline CAE-seg con backbone PanDerm Large
entrenado desde cero ($0{,}894$). La adaptación añade
$\sim 4{,}5$\,M parámetros entrenables sobre los
$\sim 227$\,M de SAM2.1-Large (proporción inferior al $2\%$),
lo que mantiene el coste computacional bajo.

La transferencia fuera de dominio es prácticamente nula: Dice
desciende sólo $0{,}002$\,pp sobre ISIC2017
($0{,}947 \to 0{,}945$) y, sobre PH2, sube hasta $0{,}960$.
Este patrón se interpreta como evidencia de que las
representaciones del encoder visual congelado (Hiera-Large)
codifican primitivas geométricas y de contorno
suficientemente generales, y que la adaptación de bajo rango sobre
el decodificador de máscaras y el \emph{promptable encoder}
captura los matices específicos del dominio dermatoscópico sin
sobreajustar a una fuente concreta de adquisición. El paso de SAM2
\emph{zero-shot} a SAM2.1-Large con LoRA introduce mejoras de
$+14{,}8$\,pp Dice sobre ISIC2017 y $+7{,}4$\,pp sobre PH2,
lo que cuantifica el aporte específico de la adaptación de
dominio.

\subsection{Sensibilidad al tokenizador en zero-shot multimodal}
\label{sec:disc-zs}

La evaluación zero-shot sobre HAM10000 (secci\'on~\ref{sec:zs-results},
tabla~\ref{tab:zs-ham}) documenta una sensibilidad acusada al
tokenizador y a la formulación de los \emph{prompts}. La AUROC
oscila entre $0{,}366$ (DermLIP v1/v2 con tokenizador PubMedBERT
y \emph{prompts} genéricos) y $0{,}854$ (DermLIP original con
tokenizador GPT-2/CLIP y \emph{prompts} A3 de Derm1M), un rango de
$48{,}8$\,pp para variaciones que no afectan a los pesos
visuales sino exclusivamente a la rama lingüística y a la
formulación textual de las clases.

El valor AUROC $<0{,}5$ con la configuración inicial es
diagnóstico: indica que el espacio textual proyectado por PubMedBERT
no está alineado con el espacio visual del encoder de DermLIP
v1/v2, lo que resulta en una correlación efectivamente invertida
respecto a la tarea de clasificación. Dos modificaciones aisladas
permiten interpretar el origen de la mejora: el cambio del conjunto
A3 manteniendo el tokenizador PubMedBERT eleva AUROC a $0{,}516$
($+15{,}0$\,pp), mientras que el cambio de tokenizador
manteniendo el conjunto A3 alcanza $0{,}854$ ($+33{,}8$\,pp
adicionales). El factor crítico es por tanto la compatibilidad del
tokenizador con el espacio de \emph{embeddings} del modelo
contrastivo. La implicación práctica es que cualquier despliegue
zero-shot debe verificar empíricamente la combinación
modelo-tokenizador-\emph{prompt} sobre un conjunto de validación
del dominio, dado que la transferencia desde la configuración
publicada por los autores no está garantizada.

La evaluación adicional sobre DermapixelAI~1.0
(secci\'on~\ref{sec:dermapixel-zs}, tabla~\ref{tab:dermapixel-zs})
extiende esta lectura hacia la dimensión idiomática. DermLIP v2
con \emph{prompts} en inglés alcanza Acc@1 $= 0{,}500$ sobre la
categoría etiológica L1, frente a $0{,}361$ con \emph{prompts}
equivalentes en castellano ($-13{,}9$\,pp); SigLIP-SO400M
presenta una caída de $-11{,}1$\,pp en la misma comparación.
BiomedCLIP mantiene Acc@1 $= 0{,}389$ en ambos idiomas, lo que
sitúa la penalización por castellano como dependiente del corpus
textual de preentrenamiento. La magnitud cuantificada motiva la
construcción de modelos vision-lenguaje contrastivos con rama
textual entrenada sobre material clínico en castellano como
línea futura inmediata.

\subsection{Brecha entre modelos específicos y LLMs multimodales}
\label{sec:disc-llm}

La comparación de paradigmas sobre HAM10000
(secci\'on~\ref{sec:llm-results}, tabla~\ref{tab:llm-comparativa})
establece un orden de magnitud entre familias que es relevante
para la selección de arquitectura en escenarios clínicos. PanDerm
Base con ajuste fino y \emph{test-time augmentation} sobre
HAM10000 alcanza $0{,}920$ accuracy y $0{,}852$ BAcc; SigLIP-Large con sondeo lineal, $0{,}900$; MedGemma~27B
con LoRA sobre HAM10000, $0{,}802$; MedGemma~27B sin
\emph{fine-tuning}, $0{,}665$; GPT-4o \emph{zero-shot},
$0{,}485$; Gemini~2.5~Pro \emph{zero-shot}, $0{,}403$. El
gradiente de rendimiento es de $43{,}5$\,pp en accuracy entre
el mejor especialista ajustado y el mejor LLM multimodal
\emph{zero-shot} ($0{,}920 - 0{,}485$).

Tres observaciones cuantitativas complementan esta brecha. Primero,
la adaptación con LoRA sobre MedGemma~27B sobre el \emph{split} de
entrenamiento de HAM10000 produce mejoras de $+13{,}7$\,pp Acc
sobre HAM10000 y $+18{,}3$\,pp Acc sobre DDI respecto a la
variante \emph{zero-shot}, sin mejora apreciable sobre
PAD-UFES-20 ($-1{,}7$\,pp). El comportamiento es asimétrico: la
adaptación captura patrones del dataset de entrenamiento sin
transferencia sistemática al resto. Segundo, GPT-4o y
GPT-4o-mini presentan patrones de error con sesgo hacia clases
minoritarias del \emph{prompt}: GPT-4o emite $284$
predicciones de dermatofibroma sobre $8$ muestras reales y
$299$ predicciones de melanoma sobre $70$ muestras, lo que
sugiere que el modelo responde al espacio léxico del \emph{prompt}
en lugar de a la evidencia visual. Tercero, BLIP-2 e InstructBLIP
con \emph{Q-Former} preentrenado producen accuracy
$\sim 0{,}02$ sobre HAM10000 y PAD-UFES-20, comportamiento
incompatible con cualquier uso clínico y atribuible a la falta de
adaptación al vocabulario dermatológico.

El coste agregado de la evaluación con APIs comerciales sobre los
tres datasets es de aproximadamente \$4{,}07. MedGemma~27B
opera localmente en BF16 sobre los $128$\,GB de memoria
unificada de la DGX Spark sin cuantización, con latencia media de
$1{,}49$ segundos por imagen. La diferencia operativa es
relevante para un despliegue hospitalario: los modelos cerrados
añaden coste recurrente y dependencia de un proveedor externo, sin
mejora apreciable de rendimiento clínico.

\subsection{Equidad por fototipo cutáneo}
\label{sec:disc-equidad}

La evaluación sobre Fitzpatrick17k completo
(secci\'on~\ref{sec:equidad-results}, tablas~\ref{tab:equidad-fototipo}
y~\ref{tab:cross-eval}) reproduce y cuantifica un patrón conocido
en la literatura: los cinco encoders contrastados presentan
\emph{gap}~I-VI positivo en la formulación de tres clases de
malignidad, con valores entre $+0{,}124$ y $+0{,}306$\,pp de
BAcc. La interpretación del \emph{gap} requiere considerar la
BAcc global de cada modelo. DermLIP v2 obtiene la mejor BAcc
global ($0{,}721$) y la mejor AUROC global ($0{,}909$), pero
su \emph{gap}~I-VI es $+0{,}281$\,pp, indicando que el
rendimiento sobre FST~VI ($0{,}434$) cae substancialmente
respecto a FST~III ($0{,}765$). PanDerm Large presenta un
\emph{gap} menor ($+0{,}217$\,pp) con BAcc global ligeramente
inferior ($0{,}699$); BiomedCLIP, el \emph{gap} mínimo
($+0{,}124$\,pp) pero la peor BAcc global ($0{,}590$).

Esta última asimetría ---\emph{gap} mínimo con BAcc global
inferior--- ilustra el carácter no decisional de la métrica de
equidad aislada. Un modelo uniformemente inferior entre subgrupos
puede presentar baja desigualdad aritmética sin ser equitativo en
sentido sustantivo, dado que su rendimiento es bajo en todos los
fototipos por igual. La interpretación operativa de las cifras
exige por tanto considerar BAcc global y \emph{gap} de forma
conjunta. El modelo con mejor compromiso entre ambas variables es
PanDerm Large, seguido de DermLIP v2 cuando se ponderan el
rendimiento sobre FST~I-IV y la AUROC global.

El experimento cruzado \emph{Light}~$\leftrightarrow$~\emph{Dark}
(tabla~\ref{tab:cross-eval}) aporta una segunda lectura: la caída
de rendimiento al entrenar sólo con fototipos claros y evaluar
sobre fototipos oscuros (L$\to$D vs L$\to$L) es máxima en DINOv2
($\Delta = -0{,}132$\,pp BAcc) y mínima en BiomedCLIP
($\Delta = -0{,}063$\,pp). DermLIP v2 mantiene la menor caída
Dark$\to$Light vs Dark$\to$Dark ($\Delta = +0{,}016$\,pp), lo que
sugiere mayor invariancia de sus representaciones al fototipo. El
patrón es coherente con la composición de su corpus de
preentrenamiento Derm1M, derivado de fuentes web abiertas con
representación variable de fototipos.

\subsection{Estimación robusta del rendimiento sobre corpus de tamaño moderado}
\label{sec:disc-cv-robustez}

La aplicación de validación cruzada de cinco \emph{folds} agrupada
por caso clínico sobre DermapixelAI~1.0
(secci\'on~\ref{sec:dermapixel-results},
tabla~\ref{tab:dermapixel-cv}) constituye un argumento
metodológico relevante para la evaluación de modelos
fundacionales sobre corpus de tamaño moderado. La diferencia entre
la cifra puntual del split fijo y la media de los cinco
\emph{folds} es sustancial sobre BAcc ($-29{,}5$\,pp en L1) y
prácticamente estable sobre AUROC ($-4{,}3$\,pp en L1), patrón
que cuantifica la sensibilidad de las métricas dependientes de
umbral al cambio de realización muestral cuando el conjunto de
test es reducido ($N = 36$). La interpretación operativa es
doble. \textit{Primero}, la cifra del split publicado es válida
como punto operativo de referencia comparable con la evaluación
\emph{zero-shot} (secci\'on~\ref{sec:dermapixel-zs}, que utiliza
el mismo test para trazabilidad cruzada), pero requiere
acompañarse de la estimación cruzada para defender el rendimiento
esperado del sistema. \textit{Segundo}, la AUROC
\emph{one-vs-rest} resulta la métrica primaria recomendable para
comparar configuraciones de \emph{transfer learning} congelado
sobre corpus clínicos de tamaño moderado, mientras que la BAcc y
la \emph{accuracy} deben reportarse acompañadas de su intervalo de
confianza y, cuando sea posible, de su estimación cruzada por
\emph{folds}. La práctica refuerza la fiabilidad del benchmark y
delimita explícitamente el alcance interpretativo de las cifras
puntuales.

\section{Comparación con la literatura}
\label{sec:lit}

\subsection{Reproducción del benchmark de PanDerm}
\label{sec:lit-panderm}

El trabajo original de Yan et~al.~\cite{panderm2025} reporta para
PanDerm Large sobre los $28$ \emph{downstream tasks} de su
benchmark interno una mejora media del $+10{,}2\%$ sobre el
estado del arte previo. Las cifras de sondeo lineal de la
tabla~\ref{tab:lp-panderm} se sitúan dentro del margen reportado
por los autores para los datasets compartidos (HAM10000,
PAD-UFES-20, Derm7pt) y reproducen cualitativamente el
comportamiento del modelo en arquitectura aarch64 con CUDA~13.0 y
BF16. La verificación numérica respecto al hardware x86\_64
canónico es por tanto satisfactoria dentro de las tolerancias
esperables por aritmética BF16 y por diferencias de versiones de
\texttt{cuDNN}, lo que documenta la portabilidad del pipeline al
chip Grace Hopper GB10.

La extrapolación de la cifra agregada del paper ($+10{,}2\%$)
a la diferencia PanDerm Large vs PanDerm Base sobre los diez
datasets externos contrastados de este trabajo
($+4{,}1$\,pp accuracy, $+9{,}0$\,pp BAcc) es coherente:
$+9{,}0$\,pp BAcc es próximo al $+10{,}2\%$ que los autores
reportan sobre su benchmark interno, lo que sugiere que la ventaja
del modelo Large es estable entre los conjuntos públicos
contrastados y los conjuntos internos del grupo de Monash.

\subsection{Segmentación binaria de lesión}
\label{sec:lit-seg}

La cifra de SAM2.1-Large con LoRA sobre ISIC2018 ($0{,}947$ Dice)
supera en $+2{,}6$\,pp el resultado publicado por
Yan et~al.~\cite{panderm2025} ($0{,}921$ Dice con $100$ épocas).
Tres factores plausibles explican la diferencia: (i)~la
arquitectura promptable de SAM2.1 con \emph{bounding box} derivada
de la máscara de \emph{ground truth} aporta un guiado espacial
fuerte que las arquitecturas no promptables no aprovechan;
(ii)~el encoder Hiera-Large preentrenado sobre $\sim 11$
millones de imágenes y mil millones de máscaras codifica
primitivas geométricas más informativas que las que el backbone
visual de PanDerm aprende sobre dermatoscopia; (iii)~la adaptación
LoRA introduce un número de parámetros suficiente para capturar las
particularidades del dominio sin sobreajuste, mientras que un
\emph{fine-tuning} completo sobre $50$ épocas podría requerir
más datos o regularización adicional para superar la cifra de
referencia. La interpretación operativa es que la segmentación
binaria de lesión sobre dermatoscopia se encuentra cerca de la
saturación cuando se combinan un encoder promptable preentrenado a
gran escala y una adaptación de bajo coste sobre ISIC2018.

\subsection{Comportamiento zero-shot}
\label{sec:lit-zs}

La AUROC máxima alcanzada en este trabajo sobre HAM10000 con
DermLIP en formulación zero-shot ($0{,}854$) coincide
cualitativamente con los rangos reportados en el paper original de
Derm1M~\cite{derm1m} para la configuración de mejor rendimiento.
La sensibilidad al tokenizador documentada en
secci\'on~\ref{sec:disc-zs} no está cuantificada de forma sistemática en
la literatura previa sobre la familia DermLIP, lo que constituye
una aportación empírica del presente trabajo. La diferencia de
$48{,}8$\,pp AUROC entre la peor y la mejor configuración
testada sobre el mismo encoder visual es de un orden de magnitud
similar al efecto de cambiar entre paradigmas (LP vs ZS), lo que
desplaza la discusión metodológica desde el modelo hacia su
configuración operativa.

\subsection{Disparidad por fototipo}
\label{sec:lit-equidad}

El sentido de la desigualdad ---peor rendimiento sobre fototipos
elevados--- coincide con los hallazgos previos sobre
Fitzpatrick17k~\cite{fitzpatrick17k} y con la motivación que dio
origen a la curación de DDI~\cite{ddi}. El presente trabajo
extiende esta línea evaluando cinco encoders bajo un mismo
protocolo, lo que permite contrastes directos. La paradoja
BiomedCLIP (\emph{gap} mínimo asociado a peor BAcc global) no está
documentada de forma explícita en la literatura sobre equidad
dermatológica y constituye una observación sobre la métrica
\emph{gap} aislada que invita a reportes complementarios: BAcc
global, AUROC global y \emph{gap} deben presentarse conjuntamente.

\section{Implicaciones operativas}
\label{sec:implicaciones}

\subsection{Arquitectura defendible para sistemas de apoyo al diagnóstico}
\label{sec:disc-arq}

La combinación de los cuatro bloques de evidencia anteriores
sugiere una arquitectura técnicamente defendible para sistemas de
apoyo al diagnóstico dermatológico construidos sobre los modelos
fundacionales disponibles. La estructura recomendable es la
siguiente: encoder PanDerm Large preentrenado, congelado por
defecto y opcionalmente ajustado bajo demanda; cabeza lineal
multi-clase con regresión logística (L-BFGS, $C=1{,}0$,
\texttt{max\_iter}~$=5\,000$) cuando el conjunto de entrenamiento
disponible sea reducido ($N_{\mathrm{train}} \lesssim 500$);
cabeza ajustada con la receta de PanDerm Large
(Mixup, CutMix, \emph{weighted random sampler}, $50$ épocas)
cuando $N_{\mathrm{train}} \gtrsim 5\,000$ y exista desbalance
acusado; segmentación con SAM2.1-Large adaptado mediante LoRA
($r = 8$) cuando la tarea requiera localización de lesión.

La rama lingüística contrastiva (DermLIP) se incorpora únicamente
con prompts validados y tokenizador compatible verificado sobre un
conjunto de validación del dominio. La rama LLM multimodal
generalista (GPT-4o, Gemini, MedGemma) no se incorpora como
clasificador principal, dado el orden de magnitud de rendimiento
inferior cuantificado en secci\'on~\ref{sec:disc-llm}; su rol natural es
complementario, como generador de razonamiento textual o
explicaciones para casos previamente clasificados por el encoder
especializado.

\subsection{Armonización taxonómica y nivel de agregación}
\label{sec:disc-armonizacion}

La heterogeneidad taxonómica entre datasets, mitigada mediante la
ontología jerárquica L1/L2/L3 descrita en secci\'on~\ref{sec:ontologia},
sigue requiriendo una decisión explícita sobre el nivel de
agregación al que se realizan los contrastes. La comparación
inter-dataset es defendible a nivel L1 (cuatro familias
etiológicas) y L2 ($43$ subcategorías clínicas) sobre el corpus
armonizado de $72\,654$ imágenes; el nivel L3 ($367$
diagnósticos individuales) conserva la cola larga propia de cada
dataset y no admite contrastes directos sin agregación
adicional. La implicación práctica para un sistema de producción es
que la jerarquía debe presentarse de forma explícita al usuario
clínico: la predicción L3 acompañada de su agregación L2 y L1
proporciona robustez frente a la incertidumbre en la cola larga sin
sacrificar especificidad cuando ésta sea fiable.

\subsection{Ventana operativa por fototipo cutáneo}
\label{sec:disc-ventana}

La ventana operativa de los modelos contrastados, definida por su
rendimiento desigual entre fototipos en el experimento de equidad
(tabla~\ref{tab:equidad-fototipo}), condiciona la generalización a
poblaciones con fototipos elevados. Los cinco encoders presentan
\emph{gap}~I-VI positivo entre $+0{,}124$ y $+0{,}306$\,pp BAcc.
La paradoja BiomedCLIP ---\emph{gap} mínimo asociado a la peor
BAcc global---  ilustra que el \emph{gap} aislado no constituye
una métrica decisional. El modelo con mejor compromiso entre BAcc
global y \emph{gap} sobre la formulación de tres clases es PanDerm
Large, seguido de DermLIP v2 cuando se prioriza la AUROC global.
La implicación clínica es que el despliegue sobre poblaciones con
representación significativa de fototipos elevados debe acompañarse
de validación específica sobre subgrupos y de mecanismos de revisión
humana cuando el modelo opere en su régimen de menor confianza.

\subsection{Ensemble de seguridad para detección de melanoma}
\label{sec:disc-ensemble}

El protocolo auxiliar de ensemble (secci\'on~\ref{sec:ensemble-results},
tabla~\ref{tab:ensemble}) documenta que la combinación OR
\emph{top-3} de PanDerm Large \emph{fine-tuned} sobre HAM10000 (M1)
y SigLIP-Large SO400M con sondeo lineal alcanza \emph{recall} de
melanoma del $100\%$ ($70$ de $70$) sobre el \emph{split}
de test de HAM10000. El análisis de errores muestra que los dos
melanomas que M1 no captura en \emph{top-3} sí son capturados por
SigLIP, lo que confirma la complementariedad de los \emph{backbones}
y motiva la elección del par. M7 (clasificador unificado sobre el
corpus armonizado de $43$ clases) no añade melanomas adicionales
sobre HAM, pero mantiene valor en escenarios multidataset por su
mayor cobertura ontológica.

Tres caveats acompañan la interpretación operativa de este
resultado: (i)~el \emph{recall} del $100\%$ es \emph{top-3}, no
\emph{top-1}, lo que requiere una interfaz que presente las tres
clases más probables al clínico; (ii)~el tamaño muestral de
melanomas es $N = 70$, lo que produce intervalos de confianza al
$95\%$ amplios y exige validación prospectiva sobre cohorte
mayor; (iii)~el coste en falsos positivos del esquema OR
\emph{top-3} es elevado ($885$ FP sobre $1\,162$ no-melanomas),
lo que se acepta en un escenario de cribado oportunista pero
descarta su uso aislado para clasificación primaria.

\section{Hallazgos no anticipados}
\label{sec:hallazgos}

Cinco observaciones del trabajo no se anticipaban a partir de la
literatura y merecen atención específica.

\paragraph{Competitividad de ConvNeXt-Large sobre HAM10000.}
ConvNeXt-Large preentrenado sobre ImageNet-21K alcanza
$0{,}967$ AUROC sobre HAM10000, superando a PanDerm Large
($0{,}954$) por $1{,}3$\,pp y aproximándose a SigLIP-Large
($0{,}971$). El resultado sugiere que para tareas de ranking
binario o multiclase sobre dermatoscopia estandarizada, las
representaciones de ImageNet-21K son altamente competitivas, y
que la ventaja del preentrenamiento de dominio aparece sólo en
material clínico no estandarizado.

\paragraph{Sondeo lineal supera al ajuste fino en DDI.}
Con $N_{\mathrm{train}} = 427$, el sondeo lineal sobre PanDerm
Large supera al ajuste fino en accuracy ($+7{,}3$\,pp), AUROC
($+14{,}5$\,pp), BAcc ($+2{,}1$\,pp) y F1 ponderada
($+5{,}6$\,pp). El comportamiento, atribuible a sobreajuste,
no se documenta de forma explícita en la literatura sobre PanDerm
y proporciona una heurística operativa para datasets pequeños.

\paragraph{Asimetría rendimiento-equidad en DermLIP v2.}
DermLIP v2 alcanza la mejor BAcc global sobre Fitzpatrick17k
($0{,}721$) y la mejor AUROC global ($0{,}909$), pero su
\emph{gap}~I-VI es $+0{,}281$\,pp, el segundo más alto entre
los cinco encoders contrastados. La combinación de rendimiento
agregado superior con desigualdad por subgrupo elevada sugiere que
el preentrenamiento sobre Derm1M produce representaciones
informativas sobre la categoría diagnóstica pero no
suficientemente invariantes al fototipo cutáneo.

\paragraph{Adaptación LoRA sobre MedGemma~27B.}
La adaptación LoRA sobre MedGemma~27B sobre el \emph{split} de
entrenamiento de HAM10000 produce mejoras de $+13{,}7$\,pp Acc
sobre HAM10000 y $+18{,}3$\,pp Acc sobre DDI respecto a la
variante \emph{zero-shot}, sin mejora apreciable sobre PAD-UFES-20
($-1{,}7$\,pp). El comportamiento es asimétrico y sugiere que
la adaptación captura patrones específicos del dataset de
entrenamiento sin transferencia sistemática a otros conjuntos. La
implicación es que el ajuste con LoRA sobre LLMs multimodales
debe reportarse por dataset, sin agregar como un único valor de
rendimiento.

\paragraph{Interpretabilidad emergente del SAE Large sobre
Derm7pt.}
La validación independiente del SAE Large
($16\,384$ \emph{features}) sobre el corpus
Derm7pt~\cite{kawahara2019} reportada en la
secci\'on~\ref{sec:sae-derm7pt} identifica \emph{features}
individuales con AUROC efectivo hasta $0{,}823$ sobre
\emph{vascular structures} y hasta $0{,}926$ sobre
vasos en horquilla ($N_{+} = 15$ positivos), sin que las
anotaciones del \emph{Seven-Point Checklist} se hayan utilizado
en el preentrenamiento del SAE ni del encoder PanDerm Large.
La implicación es que la base sparse aprendida sobre Derm1M
($507\,657$ imágenes) ha capturado representaciones alineadas
con la taxonomía dermatoscópica clásica, lo que sostiene la
línea de interpretabilidad mediante Sparse Autoencoders más allá
del régimen de los $48$ conceptos visuales de
SkinCon~\cite{skincon} sobre el que se diseñó el experimento
original del Anexo~\ref{cap:anexof}. El \emph{Concept Bottleneck
Model} con cuello de botella sobre los siete criterios cuantifica
el coste de la interpretabilidad sobre la tarea binaria
melanoma/no-melanoma ($-5{,}8$\,pp AUROC, $0{,}890 \to 0{,}832$)
y produce una jerarquía de pesos del meta-clasificador
(\emph{blue whitish veil} $+2{,}54$, \emph{regression
structures} $+2{,}07$) coherente con el peso clínico de
estos criterios en la escala diagnóstica de
referencia~\cite{kawahara2019}. La implicación operativa es que el
SAE no constituye solo un mecanismo de reducción dimensional sino
que admite interpretación concepto-a-concepto sobre el vocabulario
dermatoscópico, lo que fortalece la viabilidad del módulo de
interpretabilidad del prototipo (Anexo~\ref{cap:anexoh}) sin
requerir reentrenamiento del SAE.

\section{Limitaciones interpretativas}
\label{sec:disc-limitaciones}

La interpretación de los resultados anteriores está condicionada
por cuatro elementos que el capítulo~\ref{cap:limitaciones}
desarrolla en detalle. Primero, el solapamiento documentado entre
Dermnet, HIBA, MSKCC y el corpus de preentrenamiento Derm1M
introduce una incertidumbre acotada sobre las cifras de la
familia DermLIP en estos tres datasets; las cifras de PanDerm,
DINOv2, ConvNeXt-Large, EfficientNetV2-Large y SigLIP-Large no se
ven afectadas. Segundo, los tamaños muestrales por subgrupo
Fitzpatrick V y VI ($169$ y $73$ imágenes respectivamente
en el conjunto de test) producen intervalos de confianza amplios
para las cifras de \emph{gap}, lo que limita la cuantificación
exacta de la desigualdad por fototipo. Tercero, la sensibilidad
al \emph{prompt} en zero-shot y en la evaluación de LLMs
multimodales implica que las cifras reportadas son condicionales a
una formulación concreta, y que la generalización a otras
formulaciones requiere validación específica. Cuarto, el trabajo
evalúa rendimiento sobre fotografías procedentes de archivos de
investigación; la generalización a imagen real adquirida en
flujos de teledermatología sobre cohorte hospitalaria prospectiva
no se aborda y constituye una línea futura inmediata.

Estas limitaciones acotan el dominio de validez de las cifras
reportadas pero no invalidan los patrones cualitativos
identificados: la ventaja del preentrenamiento de dominio, la
saturación temprana del sondeo lineal, el régimen
$N_{\mathrm{train}} \lesssim 500$ para preferencia de LP, la
brecha entre especialistas y LLMs multimodales generalistas, y la
existencia de un \emph{gap} por fototipo sistemático y consistente
entre encoders.
